\section{What the evidence supports}
\paragraph{Context quality is an actionable engineering target.}
The controlled contrast tests whether preserving connected prerequisite evidence helps the same reader produce grounded decisions. The evidence/reader decomposition locates lost value: missing support cannot be repaired by a stronger reader without additional information. This supports measuring retrieval against downstream decision requirements, as well as relevance or answer recall. A deployment can retain a general company store while adding task-aware selection and explicit citations.

\paragraph{The generic task-word control bounds the specialization claim.}
An exact input match is a stronger mechanistic observation than a non-significant performance difference. When the generic query-word selector and GTM selector produce identical packets, their different labels do not create different information for the reader. We do not count aliased responses as additional experiments or claim an architectural moat from that comparison. Domain vocabulary and dependency-preserving packing are useful design choices; the experiment does not establish that a separate vertical database, proprietary graph or complex ontology is required.

\paragraph{Why the plots use policy value rather than abstention accuracy.}
In the earlier balanced eligibility benchmark, a model with no company evidence scored half the cases correctly by abstaining on negative examples. That number did not describe usable recommendations. The new primary graph awards value only for fully supported positive decisions and reports safety separately. It makes the information contribution visible while avoiding the interpretation that company-specific facts are a general language-model capability test.

\paragraph{Scope and unresolved evidence.}
All main histories are synthetic and follow a disclosed schema. We share a correct temporal/access boundary and verified relation fields, so ingestion quality and maintenance costs are largely outside the experiment. The compound stress levels combine several factors. Entity closure and query-word filtering are particular baseline implementations; GBrain is a search component here, not its complete autonomous product. The external wiki track remains small, developmental and from one company. There is no practitioner panel, live campaign, observed revenue, retention estimate or network-effect evidence. One response per condition does not measure repeated-run reliability. Prices are list-rate estimates from usage, not invoices or full infrastructure cost.

The product implication is to instrument evidence selection, source completeness and context cost before investing in more elaborate memory infrastructure. The release provides the data, reference solver, generic and task-conditioned selectors, actual model outputs and analysis needed to challenge that conclusion. Independent authoring, blinded practitioner review, raw-document ingestion and prospective business outcomes are the next validation steps.
